% ============================================================================
% Appendix A8 — Bifurcation-Interval Honesty: the Retracted beta_crit
% ============================================================================
\section{Bifurcation-Interval Honesty: the Retracted
$\beta_{\mathrm{crit}}\approx0.1015$}
\label{app:bifurcation}

Earlier drafts reported a critical coupling
$\beta_{\mathrm{crit}}\approx0.1015$ at which the equilibrium vitality
$V^{*}(\beta)$ appeared to cross the strategic threshold $\Vstrat$,
presented as a bifurcation of the baseline model. The claim was wrong
twice over---once numerically and once structurally---and this appendix
documents both failures and the corrected reporting protocol, because the
failure mode (a scan-window artifact masquerading as a critical point) is
generic enough to be worth recording.

\subsection{The structural impossibility}

By Corollary~\ref{cor:scope}, $V^{*}$ is invariant in $\beta$ on the
interior branch: at baseline it is the constant $4.7025$
($\varepsilon=0.1$) for \emph{every} $\beta$ in the structurally interior
region. A $V^{*}(\beta)$ crossing of $\Vstrat=5$ therefore cannot exist at
this calibration. Independently, Lemma~\ref{lem:trace} gives
$\det J>0$ everywhere, so no fold in $\beta$ (or any other parameter)
exists either. Whatever the legacy sweep detected, it was not a
bifurcation.

\subsection{The numerical forensics}

The legacy equilibrium scan searched for nullcline crossings in the window
$C\in[0.01,\,80]$. But the interior branch obeys
\[
C^{*}(\beta)\;=\;\frac{\beta C^{*}}{\beta}\;=\;\frac{8.008}{\beta}
\qquad(\varepsilon=0.1),
\]
which exceeds the window ceiling $C_{\max}^{\mathrm{scan}}=80$ precisely
when
\[
\beta \;<\; \frac{8.008}{80} \;=\; 0.1001 .
\]
For $\beta$ below this value the scan silently \emph{lost} the interior
equilibrium; the interpolation across the last resolved grid points then
manufactured an apparent threshold crossing at $\beta\approx0.1015$. The
``critical coupling'' was the image of the scan window, not of the model:
false precision produced by an implementation constant.

\subsection{The corrected protocol}

The companion function
\texttt{dlvt.analysis.estimate\_bifurcation\_interval} replaces the point
estimate with two honest questions: (i)~does a $V^{*}$-crossing of
$\Vstrat$ exist anywhere in the sweep, once the scan window is made large
enough ($C\le500$) to hold the low-$\beta$ branch? and (ii)~if a crossing
exists (non-baseline calibrations), what is the \emph{interval} of
numerical estimates across a grid of regularizations
$\varepsilon\in\{0.01,0.05,0.1,0.2\}$ and scan resolutions---the
methods-honest reporting granularity, in place of a four-significant-figure
point. At baseline the answers are: no crossing exists;
$V^{*}(\beta)$ is numerically constant at $4.7025$; and the conserved
product is $\beta C^{*}=8.0083$ (all three pinned by the
\texttt{test\_bifurcation\_interval\_*} group, whose diagnostic string is
required to name the scan-window artifact explicitly so the legacy figure
cannot be reproduced by accident).

\subsection{Where real $\beta$-sensitivity lives}

The retraction is scoped, not global. Outside the separable power-law
family---exponential or Hill (saturating) load/drain kernels in which the
coupling parameter also sets a saturation or curvature scale---the
$\beta$-invariance of Corollary~\ref{cor:scope} fails,
$\partial V^{*}/\partial\beta\neq0$, and genuine threshold crossings in
$\beta$ can occur (Proposition~P4, \S\ref{sec:propositions}). Establishing
the empirical kernel family is therefore a first-order measurement task,
not a modeling nicety: it decides whether scope redesign has any
equilibrium leverage at all.
